% =============================================================================
% 1.10 Workflows déterministes, architectures agentiques et HITL
% =============================================================================
\section{Workflows, architectures agentiques et supervision humaine}
\label{sec:ch1-workflows}

\subsection{Systèmes déterministes et systèmes agentiques}

La littérature récente distingue deux régimes de conception des systèmes
fondés sur les grands modèles de langage.

Dans un \btpterm{workflow déterministe}, le chemin d'exécution est connu à
l'avance. Le modèle intervient \emph{dans} une étape --- extraire, classer,
rédiger --- mais ne décide pas de l'enchaînement des étapes. Le code
orchestre, filtre, valide et interrompt.

Dans un \btpterm{système agentique}, le modèle décide de la prochaine
action, peut boucler, reformuler, changer d'outil et d'objectif. Cette
flexibilité est précieuse lorsque le chemin n'est pas prédictible. Elle a
un coût~: l'indétermination du parcours, la difficulté d'audit et le risque
d'une action non autorisée.

Anthropic (2024) souligne que les systèmes efficaces sont souvent construits
à partir de composants simples et composables, et que l'ajout d'agents doit
être justifié par la complexité réelle du problème, non par une préférence
esthétique pour l'autonomie. Wang et coll.\ (2023) recensent par ailleurs
les défaillances typiques des agents autonomes~: boucles infinies, dérive
d'objectif, appels d'outils injustifiés.

\begin{btptable}{Comparaison des deux régimes}
\label{tab:workflow-vs-agent}
\begin{tabularx}{\linewidth}{@{}L{3.6cm} Y Y@{}}
\toprule
\rowcolor{btppale}
\btpth{Critère} & \btpth{Workflow déterministe} & \btpth{Agent autonome} \\
\midrule
Chemin d'exécution & Connu, reproductible. & Décidé à l'exécution. \\
Traçabilité & Chaque étape est journalisable. & Le journal est un récit, pas un contrat. \\
Contrôle d'erreur & Validations déterministes entre étapes. & Dépend de l'auto-critique du modèle. \\
Coût d'inférence & Borné par le graphe d'étapes. & Potentiellement non borné. \\
Pertinence pour les AO & Élevée~: le processus métier est connu. & Ponctuelle~: questions multi-sauts. \\
\bottomrule
\end{tabularx}
\end{btptable}

Pour les appels d'offres, le processus métier est connu~: déposer, analyser,
extraire, rechercher, rédiger une section, faire valider, exporter. Ce
n'est pas un espace d'exploration. Une approche workflow, dans laquelle
l'intelligence est \emph{locale} à chaque étape, est donc le régime par
défaut. Un n\oe ud agentique peut, plus tard, être encapsulé dans une
étape précise --- par exemple une recherche multi-sauts --- sans changer le
reste du processus.

\subsection{Supervision humaine dans la boucle}

La supervision humaine (\anglais{Human-In-The-Loop}, HITL) désigne
l'insertion d'un opérateur dans le cycle de production, à des points
d'interruption définis. Elle n'est pas un aveu d'échec de l'automatisation~:
elle est la traduction opérationnelle du caractère engageant de l'offre.

Trois modalités se distinguent~:

\begin{description}
  \item[Approbation.] L'opérateur accepte la section telle quelle.
  \item[Correction.] L'opérateur modifie le texte~; la version corrigée
        devient la référence, et le système en conserve la trace.
  \item[Relance.] L'opérateur rejette la section et redéclenche la
        génération, éventuellement avec une consigne supplémentaire.
\end{description}

\begin{btpbox}[Ce que la supervision humaine garantit]
Elle garantit trois propriétés que aucun modèle ne peut certifier seul~:
la \btpterm{responsabilité} (un humain signe l'offre), la
\btpterm{conformité} (un expert vérifie que la section répond à
l'exigence) et l'\btpterm{apprentissage organisationnel} (les corrections
alimentent le patrimoine pour les dossiers suivants).
\end{btpbox}

\subsection{Architectures pour la génération de documents}

La génération d'un document long par un modèle de langage se heurte à une
limite désormais bien documentée~: la dégradation de l'attention au milieu
des contextes longs (Liu et coll., 2023). Plus le document à produire est
étendu, plus une génération globale perd en fidélité, en structure et en
couverture des exigences.

Trois stratégies s'opposent~:

\begin{btptable}{Stratégies de génération de documents longs}
\label{tab:strategies-generation}
\begin{tabularx}{\linewidth}{@{}L{3.6cm} Y Y@{}}
\toprule
\rowcolor{btppale}
\btpth{Stratégie} & \btpth{Avantage} & \btpth{Limite} \\
\midrule
Génération globale & Simple, un seul appel. & Perte de structure, citations diluées, pas de versions intermédiaires. \\
Génération par section & Contrôle local, citations ciblées, versions. & Exige un plan préalable et une orchestration. \\
Génération par section avec revue & Idem, plus responsabilité humaine. & Temps de cycle plus long. \\
\bottomrule
\end{tabularx}
\end{btptable}

Pour un dossier d'offre, dont le plan est largement imposé par le dossier
entrant, la génération par section n'est pas un raffinement~: elle est la
seule stratégie qui rende possibles la citation, la revue et le versionnage.
Nous en ferons le c\oe ur de la méthodologie au chapitre~\ref{chap:methodologie}.
